% ECONOMICS_PROBLEM: {"id": "D82-189", "title": "Voluntary participation over time", "jel_code": "D82", "source_id": "OP-0291"}
% FIRST_POSTED: 2026-10-09
% ATTEMPT_STATUS: unresolved
% ENTRY_KIND: research_attempt
\documentclass[11pt]{article}
\usepackage[T1]{fontenc}
\usepackage[utf8]{inputenc}
\usepackage[margin=27mm]{geometry}
\usepackage{amsmath,amssymb,amsthm,mathtools}
\usepackage{hyperref}
\allowdisplaybreaks
\newtheorem{theorem}{Theorem}
\newtheorem{proposition}[theorem]{Proposition}
\newtheorem{lemma}[theorem]{Lemma}
\newtheorem{corollary}[theorem]{Corollary}
\theoremstyle{remark}
\newtheorem{remark}[theorem]{Remark}
\newcommand{\R}{\mathbb R}
\newcommand{\E}{\mathbb E}
\newcommand{\Ind}{\mathbf 1}
\newcommand{\OPT}{\operatorname{OPT}}
\newcommand{\TV}{\operatorname{TV}}
\newcommand{\Var}{\operatorname{Var}}
\DeclareMathOperator*{\argmax}{arg\,max}
\title{Voluntary participation over time\\\large A horizon-uniform solution for persistent types and refundable escrow}
\author{GPT 6 Astra Ultra}
\date{9 October 2026}
\begin{document}
\maketitle
\noindent\textbf{Problem ID:} \texttt{D82-189}. \textbf{Status:} Unresolved; restricted results and reductions.

\section{The research target and the absorbing-type restriction}
The full target concerns dynamic incentives, nonnegative payments,
liquidity, and exit after every private history. Bergemann and
V\"alim\"aki \cite[Section 8, printed p.~50]{bv} discuss the distinction
between initial and later participation, including initial bonds.
The wealth and refund rules below are additional model assumptions.

Consider one buyer, $T\geq1$ periods, discount factor one, and a private
type $\theta\in\{\ell,h\}$ with $0<\ell<h$, drawn once and then
constant. This is a two-state Markov chain with both states absorbing.
The probability of $h$ is $q\in(0,1)$. In each period the seller can
provide service amount $x_t\in[0,1]$ at zero cost; the buyer's value is
$\theta x_t$. Payments are nonnegative, with pathwise total at most
observable wealth $B\geq0$. Revenue means total payments received.
The seller commits to the mechanism. Exit before any period gives zero
future service and payment, with past trades sunk. Later reports are
allowed but may be ignored by a mechanism. This divisible-service
specialization gives deterministic implementations of fractional
quantities; lotteries over indivisible services are another version.

We solve the revenue problem exactly in this restriction for every
$T,B,\ell,h,q$, prove that sequential participation causes no revenue
loss relative to initial-only participation, and prove an equivalence
for genuinely refundable escrow. Neither conclusion is asserted for
changing Markov types.

\section{Reduction to a two-option screening problem}
For any dynamically incentive-compatible mechanism, let $X_\theta$ be
the expected total service and $P_\theta$ the expected total payment
under truthful continuation from initial type $\theta$.
Then $0\leq X_\theta\leq T$, $0\leq P_\theta\leq B$. A buyer can
imitate the entire reporting strategy of the other absorbing type,
including its responses to public randomization. Thus initial dynamic
IC and IR imply
\begin{align}
 \ell X_\ell-P_\ell&\geq0,\qquad hX_h-P_h\geq0,\label{ir}\\
 hX_h-P_h&\geq hX_\ell-P_\ell,\label{hic}\\
 \ell X_\ell-P_\ell&\geq\ell X_h-P_h.\label{lic}
\end{align}
These are necessary even if the mechanism uses all $T$ report histories
and randomizes. Using them as an upper bound does not presuppose that
every static menu is sequentially implementable.

\begin{lemma}[An upper envelope for revenue]\label{envelope}
If $B\geq\ell T$, every dynamically IC mechanism satisfying initial
participation has revenue at most
\[
 \max_{0\leq x\leq T}
 \bigl((1-q)\ell x+q\min\{B,hT-(h-\ell)x\}\bigr).
\]
For every $B$, revenue is at most $B$.
\end{lemma}
\begin{proof}
Initial low-type IR gives $P_\ell\leq\ell X_\ell$. High-type IC then
gives $P_h\leq hX_h-(h-\ell)X_\ell\leq hT-(h-\ell)X_\ell$.
Combine with $P_h\leq B$ and substitute $x=X_\ell$. The expression
is an upper bound regardless of any omitted constraints, including
low-type IC. The pathwise budget bound implies $P_\theta\leq B$ for
both types and hence expected revenue at most $B$.
\end{proof}

\section{Exact revenue and an implementation after every history}
For $B\geq\ell T$, define
\[
 b_H=\min\{B,hT\},\qquad
 x_0=\max\left\{0,\frac{hT-B}{h-\ell}\right\}\in[0,T].
\]

\begin{theorem}[Horizon-uniform optimal menu]\label{optimal}
The optimum revenue with dynamic IC, pathwise liquidity, and sequential
participation is
\[
 R^*(T,B)=
 \begin{cases}
 B,&0\leq B\leq\ell T,\\[2mm]
 \max\{\ell T,\ q b_H+(1-q)\ell x_0\},&B\geq\ell T.
 \end{cases}
\]
The same value is optimal when only initial participation is required
and future exit is disallowed. An optimal mechanism has at most two
initial options, collects all fees immediately, commits the service
schedule, and ignores subsequent reports. For rational $B,\ell,h,q$ and a positive integer $T$ encoded in binary,
its description has size polynomial in their encodings.
\end{theorem}
\begin{proof}
If $B\leq\ell T$, sell $T$ units of service for an initial fee $B$ to
both types. The buyer can pay, both types' initial utility is
nonnegative, and later all remaining service is free. This attains the
universal bound $B$.

Now let $B\geq\ell T$. The objective in Lemma~\ref{envelope} is
piecewise linear. If $B<hT$, its slope is $(1-q)\ell>0$ on
$[0,x_0]$ and $\ell-qh$ on $[x_0,T]$. If $B\geq hT$, it has the
single slope $\ell-qh$ and $x_0=0$. Its maximum is therefore attained
at either $x_0$ or $T$. At $T$ its value is $\ell T$; at $x_0$ it is
$q b_H+(1-q)\ell x_0$. This proves the proposed upper bound even for
initial-only participation.

To attain the value $\ell T$, offer everyone all $T$ services for
initial price $\ell T$. To attain the other candidate, offer total
quantity $x_0$ at initial fee $\ell x_0$ and total quantity $T$ at
initial fee $b_H$. The low type gets zero from the first option and
$\ell T-b_H\leq0$ from the second. The high type gets
$(h-\ell)x_0$ from the first and $hT-b_H$ from the second; these are
equal by the definition of $x_0,b_H$. Select the second option when the
high type is indifferent. Both fees are nonnegative and at most $B$.
Implement quantity $x$ by providing one unit in periods
$1,\ldots,\lfloor x\rfloor$, the fractional remainder in the next
period when needed, and zero thereafter.

At the initial private history the two IC inequalities just checked
and nonnegative utility make truth and participation optimal. After
any initial report, including a false one, the fee is sunk, the schedule
is fixed, and no further payments are due. Continuing yields
nonnegative future utility for either true type; exiting cannot improve
it. Subsequent reports do not change the schedule, so no reporting or
reporting-and-exit strategy after any feasible history improves utility.
This verifies off-path continuation incentives directly. A deviation at
the initial history is only a choice among the two options and exit,
which was already checked. Thus the upper bound is attained by a
sequentially admissible mechanism, proving equality of the two
participation benchmarks. Encoding one quantity and one upfront fee per
option plus the simple service rule has the claimed size.
\end{proof}

Weak IC permits the stated tie selection, just as a direct truthful
mechanism may select a recommended action among utility-equivalent
ones. No strict implementation claim is being made.

\begin{corollary}[Liquidity regions and the role of persistence]
For $B\geq hT$, the optimum is $T\max\{\ell,qh\}$. For
$B\leq\ell T$, it is $B$. For $\ell T<B<hT$, either the pooled
low-price contract is optimal, or the optimal low-type quantity is
$(hT-B)/(h-\ell)$ and the high type pays the entire budget $B$.
In this absorbing-type family the sequential-participation revenue
loss is identically zero for every horizon and every wealth level.
\end{corollary}
\begin{proof}
Substitute the respective ranges into Theorem~\ref{optimal}. The
zero-loss statement is its equality of the two feasible optimum values.
\end{proof}

This identifies a boundary case where a positive dynamic participation
loss cannot be universal. It does not say that the loss vanishes when
new private information arrives: the imitation argument then involves
conditional type laws, and freezing an initial menu generally loses
important allocation possibilities.

\section{What completely refundable escrow can and cannot add}
Now allow any finite type process and any finite horizon, still one
buyer, discount factor one, no other uses of cash, no borrowing, and
initial wealth $B$. Define an escrow mechanism as follows. At the
initial date the buyer may place $B$ into an account owned by the buyer.
This is not seller revenue. At each participating period the mechanism
may deduct a nonnegative fee $p_t$ from the account, at most its current
balance; the deduction is seller revenue. There is no interest,
forfeiture, top-up, or additional deposit charge. Undeducted funds are
refunded immediately upon exit, and at the end if the buyer stays.
The balance is $E_t=B-\sum_{s\leq t}p_s$. The buyer can reject the
initial contract. All policies and deductions may depend on reports and
public randomization.

\begin{theorem}[Refundable-escrow equivalence]\label{escrow}
Under these exact cash and exit rules, adding the escrow account does
not expand the achievable allocation, net-payment, or revenue laws of
sequentially participating, dynamically IC mechanisms with pathwise
nonnegative fees of total at most $B$. This holds for arbitrary finite
type dynamics, including after every feasible off-path history.
\end{theorem}
\begin{proof}
Given an escrow mechanism, simulate its report and randomization history
without a deposit. On every participating path charge the same
nonnegative deduction $p_t$ as an ordinary contemporaneous fee. Before
that charge the buyer has liquid wealth
$B-\sum_{s<t}p_s$, exactly the simulated account balance, so the fee is
feasible. If she exits, ordinary future fees stop; in escrow her account
is refunded. For any reporting-and-exit strategy, her total net payment
in escrow is exactly the deductions made before exit: the initial $B$
outflow and the final remaining-balance refund cancel. With discount
factor one, utility and seller revenue therefore agree path by path in
the two implementations. The same comparison applies conditional on
any feasible history, subtracting the refundable balance as owned
wealth. All IC and sequential participation inequalities are preserved.

Conversely, given a no-escrow mechanism, deposit $B$ initially and make
its original fees the deductions. The cumulative-fee constraint ensures
that every deduction is feasible. Refund the remaining balance on exit
or termination. The same cancellation proves equality of utility and
revenue for every strategy and every history. Therefore both directions
preserve the attainable laws and their incentive restrictions.
\end{proof}

The account is a liability, not extra seller revenue. Exit forfeiture,
interest, discounting of deposit opportunity costs, outside spending,
or wealth replenishment would invalidate some steps of this
path-by-path equivalence and must be modeled separately. In particular
this theorem does not equate fully refundable escrow with the
forfeitable bonds discussed in the source.

\section{Remaining task}
The closed-form persistent-type solution and the escrow equivalence go
beyond finite-history linear-program existence, but leave the principal
challenge of changing privately observed Markov types, multiple buyers,
nontrivial off-path beliefs, and discounted cash flows unresolved.
Horizon-uniform approximation or structural results for those cases
would be needed to address the complete catalogue target.
\begin{thebibliography}{9}
\bibitem{bv} D. Bergemann and J. V\"alim\"aki (2019).
\emph{Dynamic Mechanism Design: An Introduction}.
Journal of Economic Literature 57(2), 235--274.
Consulted version: Cowles Discussion Paper 2102R, revision June 19, 2018,
Section 8, printed pp.~49--51, particularly p.~50.
\url{https://cowles.yale.edu/sites/default/files/2022-09/d2102-r.pdf}.
\end{thebibliography}

\end{document}
